% Part: turing-machines
% Chapter: machines-computations
% Section: unary-numbers

\documentclass[../../../include/open-logic-section]{subfiles}

\begin{document}

\olfileid{tur}{mac}{una}
\olsection{اعداد دی اک نشانی نمائندگی}

\begin{explain}
ٹیورنگ مشیناں اپنی پٹی
اُتے لکھے نشاناں دے
سلسلّیاں اُتے کم کردیاں
نیں۔ مشین دے ورتے
حرفاں دے مجموعے دے
مطابق ایہ سلسلے کئی
طرح دے انپٹ تے آؤٹ پٹ
دکھا سکدے نیں۔ خاص طور
اُتے، \emph{حسابی}
فنکشن، یعنی قدرتی
اعداد دے فنکشن، دا
حساب کرن والیاں ٹیورنگ
مشیناں وچ ساڈی دلچسپی
اے۔ مثبت صحیح اعداد
نوں دکھان دا اک سادہ
طریقہ ایہ اے کہ ہر عدد
نوں اکّو نشان~$\TMstroke$
دی لڑی نال کوڈ کریے۔
جے $n \in \Nat$ ہووے
تاں $\TMstroke^n$، $n = 0$
ہون اُتے خالی سلسلہ
اے، تے باقی صورت وچ
ٹھیک $n$ عدد~$\TMstroke$
نشاناں دا سلسلہ اے۔
\end{explain}

\begin{defn}[حسابی کارروائی]
ٹیورنگ مشین~$M$ فنکشن
$f\colon \Nat^k \to \Nat$
دا \emph{حساب کردی اے}
اُدوں تے صرف اُدوں جدوں
$M$ ہِٹھلے انپٹ اُتے
رُکے:
\[
\TMstroke^{n_1} \TMblank \TMstroke^{n_2} \TMblank \dots \TMblank \TMstroke^{n_k}
\]
تے آؤٹ پٹ
$\TMstroke^{f(n_1, \dots, n_k)}$
ہووے۔
\end{defn}

\begin{prob}
ایہدی تعریف دیو کہ
ٹیورنگ مشین~$M$ فنکشن
$f\colon \Nat^k \to \Nat^m$
دا حساب کدوں کردی اے۔
\end{prob}

\begin{ex}\ollabel{ex:adder}
\emph{جمع:}
آؤ فنکشن $f(n,m) = n + m$
دا حساب کرن والی مشین
بنائیے۔ مشین دی پٹی
اُتے شروع وچ $\TMstroke$
نشاناں دے دو ٹکڑے
ہون، جنہاں دیاں لمبائیاں
$n$ تے~$m$ ہون، تے
رُکّن ویلے $n+m$~$\TMstroke$
نشاناں دا اک ٹکڑا رہ
جاوے۔ انپٹ دے ایہہ دو
$\TMstroke$ ٹکڑے اک~$\TMblank$
نال وکھ کیتے ہوندے
نیں۔ ایس لئی اک طریقہ
ایہ اے کہ $\TMblank$
والے خانے اُتے اک لکیر
لکھی جاوے تے آخری~$\TMstroke$
مٹا دتا جاوے۔
\begin{figure}\[
\begin{tikzpicture}[->,>=stealth',shorten >=1pt,auto,node distance=2.8cm,
                    semithick]
  \tikzstyle{every state}=[fill=none,draw=black,text=black]

  \node[initial,state]         (A)              {$q_0$};
  \node[state]         (B) [right of=A] {$q_1$};
  \node[state]         (C) [right of=B] {$q_2$};

  \path (A) edge node {\TMtrans{\TMblank}{\TMstroke}{\TMstay}} (B)
           edge [loop above] node {\TMtrans{\TMstroke}{\TMstroke}{\TMright}} (B)
        (B) edge [loop above] node {\TMtrans{\TMstroke}{\TMstroke}{\TMright}} (B)
            edge node {\TMtrans{\TMblank}{\TMblank}{\TMleft}} (C)
        (C) edge [loop above] node {\TMtrans{\TMstroke}{\TMblank}{\TMstay}} (C);
\end{tikzpicture}
\]\caption{$f(x,y) = x+y$ دا حساب کرن والی مشین}
\ollabel{fig:adder}
\end{figure}
\end{ex}

\begin{prob}
\olref[tur][mac][una]{ex:adder}
وچلی مشین دیاں مکمل
حالتاں دا سلسلہ
انپٹ~$\tuple{3,2}$ لئی
قدم بہ قدم لکھو۔ جے
مشین $0+0$ دا حساب
کرے تاں کیہ ہوندا اے؟
\end{prob}

\begin{explain}
\olref[rep]{ex:doubler}
وچ اسی اک ٹیورنگ مشین
دی مثال دتی سی جیڑی
انپٹ وجوں~$\TMstroke$
نشاناں دی لڑی لیندی
اے تے پٹی اُتے دوگنے
~$\TMstroke$ نشاناں دی
لڑی نال رُکدی اے:
دوگنا کرن والی مشین۔
پر آؤٹ پٹ وچ
$\TMblank$ نشان دوگنے
~$\TMstroke$ ٹکڑے دے
کھبے پاسے رہ جاندے
نیں۔ ایس لئی ایہ مشین
فنکشن $f(x) = 2x$ دا حساب
نہیں کردی، بھاویں تسی شاید
ایہ سمجھیا ہووے کہ کردی اے۔ اسی ایس
مسئلے نوں ٹھیک کرن دے
دو طریقے بیان کراں گے۔
\end{explain}

\begin{ex}
\olref{fig:doubler-disc}
وچلی مشین فنکشن $f(x) =
2x$ دا حساب کردی اے۔
\olref[rep]{ex:doubler}
وچلی مشین ہر انپٹ
$\TMstroke$ مٹا کے
دور سجے پاسے دو
$\TMstroke$ لکھدی اے؛
ایس دے برعکس ایہ
مشین انپٹ دی ہر
~$\TMstroke$ لئی
سجّے پاسے صرف اک
~$\TMstroke$ ودھاؤندی
اے۔ ایہنوں انپٹ
دے ختم ہون دی جگہ
یاد رکھنی پیندی اے،
ایس لئی انپٹ تے
نواں لکھیاں لکیراں
دے وچ اک~$\TMblank$
چھڈدی اے، جیہنوں
آخر وچ~$\TMstroke$
نال بھر دیندی اے۔
نالے انپٹ وچ اپنی
موجودہ جگہ یاد
رکھن لئی اسی انپٹ
دے ٹکڑے دا اک
~$\TMstroke$ عارضی
طور اُتے~$\TMblank$
نال بدل دیندے آں۔
  \begin{figure}
\[
\begin{tikzpicture}[->,>=stealth',shorten >=1pt,auto,node distance=2.8cm,
                    semithick]
  \tikzstyle{every state}=[fill=none,draw=black,text=black]

  \node[initial,state] (0)              {$q_0$};
  \node[state]         (1) [above of=0] {$q_1$};
  \node[state]         (2) [above of=1] {$q_2$};
  \node[state]         (3) [right of=2] {$q_3$};
  \node[state]         (4) [below of=3] {$q_4$};
  \node[state]         (5) [below of=4] {$q_5$};
  \node[state]         (6) [above right of=4] {$q_6$};
  \node[state]         (7) [below of=6] {$q_7$};
  \node[state]         (8) [below of=7] {$q_8$};

  \path (0) edge node {\TMtrans{\TMstroke}{\TMblank}{\TMright}} (1)
%    (0) edge node {\TMtrans{\TMblank}{\TMblank}{\TMstay}} (h)
    (1) edge [loop left] node {\TMtrans{\TMstroke}{\TMstroke}{\TMright}} (1)
      edge node {\TMtrans{\TMblank}{\TMblank}{\TMright}} (2)
    (2) edge [loop above] node {\TMtrans{\TMstroke}{\TMstroke}{\TMright}} (2)
        edge node {\TMtrans{\TMblank}{\TMstroke}{\TMleft}} (3)
    (3) edge [loop above] node {\TMtrans{\TMstroke}{\TMstroke}{\TMleft}} (3)
        edge node[left] {\TMtrans{\TMblank}{\TMblank}{\TMleft}} (4)
    (4) edge        node[left] {\TMtrans{\TMstroke}{\TMstroke}{\TMleft}} (5)
    (5) edge [loop below] node {\TMtrans{\TMstroke}{\TMstroke}{\TMleft}} (5)
        edge              node {\TMtrans{\TMblank}{\TMstroke}{\TMright}} (0)
    (4) edge      node[sloped] {\TMtrans{\TMblank}{\TMstroke}{\TMright}} (6)
    (6) edge              node {\TMtrans{\TMblank}{\TMstroke}{\TMright}} (7)
    (7) edge [loop right] node {\TMtrans{\TMstroke}{\TMstroke}{\TMright}} (7)
        edge              node {\TMtrans{\TMblank}{\TMblank}{\TMleft}} (8)
    (8) edge [loop right] node {\TMtrans{\TMstroke}{\TMblank}{\TMstay}} (8);
    \end{tikzpicture}
\]
\caption{$f(x) = 2x$ دا حساب کرن والی مشین}
\ollabel{fig:doubler-disc}
\end{figure}
\end{ex}

\begin{ex}\ollabel{ex:mover}
$f(x) = 2x$ دا حساب
کرن دا دوجا طریقہ
ایہ اے کہ اصل
دوگنا کرن والی مشین
نوں برقرار رکھیا جاوے،
پر آخر وچ ہور حالتاں
تے حکم شامل کیتے
جان جیڑے دوگنی لکیراں
دا ٹکڑا پٹی دے بالکل
کھبے پاسے لے جان۔
\olref{fig:mover} وچلی
مشین صرف ایہ آخری
کم کردی اے۔ جے شروع
وچ پٹی اُتے پہلے
$\TMblank$ نشاناں دا
ٹکڑا، فیر~$\TMstroke$
نشاناں دا ٹکڑا ہووے
(تے سر~$\TMblank$
دے ٹکڑے وچ کسے وی
جگہ ہووے)، تاں ایہ
مشین~$\TMstroke$
نشاناں نوں اک اک کر
کے مٹا کے پٹی دے
شروع وچ لکھدی اے۔
کم پورا ہون دا پتا
لاؤن لئی ایہ پہلے
$\TMstroke$ ٹکڑے دے
آخر اُتے اک~$\TMendtape$
نشان لاؤندی اے،
جیہڑا آخر وچ مٹا
دتا جاندا اے۔ اسی
حالتاں نوں نمبر دینا~$q_6$
توں شروع کیتا اے
تاں جو ایہناں نوں
دوگنا کرن والی مشین
نال جوڑ سکئیے۔ صرف
اک ہور حکم
$\delta(q_0,
\TMblank) = \tuple{q_6,\TMblank,\TMstay}$
دی لوڑ اے؛ یعنی
~$q_0$ توں~$q_6$
تک اک تیر جس اُتے
~$\TMtrans{\TMblank}{\TMblank}{\TMstay}$
لکھیا ہووے۔ (اک
باریک مسئلہ اے:
نتیجے والی مشین
انپٹ~$x=0$ لئی
کم نہیں کردی۔ ایہنوں
اسی مشق لئی چھڈدے
آں۔)
\begin{figure}
  \[
  \begin{tikzpicture}[->,>=stealth',shorten >=1pt,auto,node distance=2.8cm,
                      semithick]
    \tikzstyle{every state}=[fill=none,draw=black,text=black]
    \node[initial,state] (6)              {$q_6$};
    \node[state]         (7) [right of=6] {$q_7$};
    \node[state]         (8) [right of=7] {$q_8$};
    \node[state]         (9) [below of=8] {$q_9$};
    \node[state]         (10) [left of=9] {$q_{10}$};
    \node[state]         (11) [left of=10]  {$q_{11}$};
    \node[state]         (12) [below of=11]  {$q_{12}$};
    \node[state]         (13) [right of=12] {$q_{13}$};
    \node[state]         (14) [right of=13] {$q_{14}$};
    \path
    (6)  edge node {\TMtrans{\TMblank}{\TMblank}{\TMright}} (7)
    (7)  edge [loop above] node {\TMtrans{\TMblank}{\TMblank}{\TMright}} (7)
         edge node {\TMtrans{\TMstroke}{\TMstroke}{\TMright}} (8)
    (8)  edge [loop above] node {\TMtrans{\TMstroke}{\TMstroke}{\TMright}} (8)
         edge node {\TMtrans{\TMblank}{\TMendtape}{\TMleft}} (9)
    (9)  edge [loop right] node {\TMtrans{\TMstroke}{\TMstroke}{\TMleft}} (9)
         edge node {\TMtrans{\TMblank}{\TMblank}{\TMright}} (10)
    (10) edge node[above] {\TMtrans{\TMstroke}{\TMblank}{\TMleft}} (11)
    (11) edge [loop above] node {\TMtrans{\TMblank}{\TMblank}{\TMleft}} (11)
         edge node[left] {\begin{tabular}{@{}l@{}}
          \TMtrans{\TMendtape}{\TMendtape}{\TMright}\\
          \TMtrans{\TMstroke}{\TMstroke}{\TMright}
         \end{tabular}} (12)
    (12) edge node[] {\TMtrans{\TMblank}{\TMstroke}{\TMright}} (13)
    (13) edge [loop below] node {\TMtrans{\TMblank}{\TMblank}{\TMright}} (13)
         edge node[sloped] {\TMtrans{\TMstroke}{\TMblank}{\TMleft}} (11)
    (13) edge node {\TMtrans{\TMendtape}{\TMblank}{\TMstay}} (14);
  \end{tikzpicture}
  \]
  \caption{$\TMstroke$ نشاناں دا ٹکڑا کھبے پاسے لے جانا}
  \ollabel{fig:mover}
\end{figure}
\end{ex}

\begin{prob}
\olref[tur][mac][una]{ex:mover}
وچ اسی اوہ مشین بیان کیتی اے
جیڑی
\olref[tur][mac][una]{fig:doubler-disc}
وچلی دوگنا کرن والی
مشین تے
\olref[tur][mac][una]{fig:mover}
وچلی ٹکڑا سرکاؤن
والی مشین نوں
ملا کے بندی اے۔
جے ایہ جوڑی ہوئی
مشین انپٹ~$x=0$
اُتے، یعنی خالی پٹی
اُتے، شروع کیتی
جاوے تاں کیہ
ہوندا اے؟ مشین
نوں کِویں بدلو گے
تاں جو ایس صورت
وچ اوہ آؤٹ پٹ~$2x=0$
نال رُکے؟ (ایہ
کم اک حالت تے
اک حالتی تبدیلی
ودھا کے ہو جانا
چاہیدا اے۔)
\end{prob}

\begin{prob}
\emph{تفریق:} اک
ٹیورنگ مشین بناؤ
جیڑی لکیراں دیاں دو
غیر خالی لڑیاں
انپٹ وجوں لے،
جنہاں دیاں لمبائیاں
$n$ تے~$m$ ہون،
جتھے $n >
m$، تے فنکشن
$f(n,m) = n - m$
دا حساب کرے۔
\end{prob}

\begin{prob}
\emph{برابری:} اک
ٹیورنگ مشین بناؤ
جیڑی ہِٹھلے فنکشن
دا حساب کرے:
\[
\fn{equality}(n,m) = 
\begin{cases}
  \text{1} & \text{جے~$n = m$} \\
  \text{0} & \text{جے~$n \neq m$}
\end{cases}
\]
جتھے~$n$ تے~$m \in \PosInt$
نیں۔
\end{prob}

\begin{prob}
فنکشن $\min(x,y)$
دا حساب کرن لئی
ٹیورنگ مشین بناؤ،
جتھے $x$ تے $y$
مثبت صحیح اعداد
نیں؛ پٹی اُتے اوہ
$\TMstroke$ نشاناں
دیاں لڑیاں نال
دکھائے جاندے نیں
جنہاں دے وچ
اک $\TMblank$ ہوندا
اے۔ مشین دے
حرفاں دے مجموعے
وچ ہور نشان وی
ورت سکدے او۔

فنکشن $\min$
اپنے آرگومنٹاں وچوں
سب توں چھوٹی
قدر چُنندا اے؛
ایس لئی
$\min(3,5)=3$،
$\min(20,16)=16$،
$\min(4,4)=4$
تے ایس طرح اگے۔
\end{prob}

\begin{defn}
ٹیورنگ مشین~$M$
جزوی فنکشن
$f\colon \Nat^k
  \pto \Nat$ دا
حساب کردی اے
اُدوں تے صرف
اُدوں جدوں:
  \begin{enumerate}
    \item $M$ انپٹ
    $\TMstroke^{n_1}\concat\TMblank\concat
    \dots \concat\TMblank\concat\TMstroke^{n_k}$
    اُتے آؤٹ پٹ
    $\TMstroke^{m}$
    نال رُکدی اے، جے
    $f(n_1, \dots, n_k) = m$
    ہووے۔
    \item $M$ یا تاں بالکل
    نہیں رُکدی،
    یا ایہو جیہے آؤٹ پٹ
    نال رُکدی اے
    جیڑا صرف
    ~$\TMstroke$
    نشاناں دا اک
    ٹکڑا نہیں، جے
    $f(n_1, \dots, n_k)$
    غیر تعریف شدہ
    ہووے۔
  \end{enumerate}
\end{defn}

\end{document}
